\begin{textsample}{POS Dim 8 – global\_north\_2023\_12 – Score 4.00 – global\_north\_2023\_12/103957056.txt}  \label{ex:f8_pos_275}
Cumbrian resident sheds light on Palestinian struggles First-hand account of Palestinian struggles at upcoming Carlisle gig Nicholas Robson pictured during his time living in Jerusalem ( Image : Supplied ) A Cumbrian resident will be sharing his first-hand perspective on Palestinian struggles at an upcoming \textbf{fundraising} gig in Carlisle.
Nicholas Hope lived in Jerusalem ' s Ramallah back in 2010 where he worked for the European Union in their police and justice mission in Palestine."I ' ve not visited Gaza, but I ' ve travelled all over the West Bank,"he said."Back then, you could see the infrastructure and how the settlers were being prioritised and protected.
They had all the highways and infrastructure - which Palestinians couldn't use."You got a palpable sense that things were getting worse for ordinary Palestinians,"he said.
During his tenure, the EU mission aimed to facilitate a two-state solution, but Mr Hope said he observed a troubling \textbf{trend}."The EU ' s mission at the time was trying to make a two-state solution happen @ @ @ @ @ @ @ @ @ @ wrong direction."I returned to the region in 2018, whilst working with the University of Cumbria teaching law."I could see for myself then that things had really gotten bad."Brand new settlements were springing up all over the West Bank, it had me wondering how the two-state solution was ever going to work."The idea you could divide two populations to make a coherent state for Palestine just seemed absurd,"he said.
Speaking about life in Gaza, Mr Hope described it as"intolerable"with restrictions on travel, family visits, and educational opportunities for Palestinians.
He also criticised what he referred to as the"cutting the grass approach"which he claimed has been adopted in Israel, with periodic bombings and military actions against the people of Gaza every few years.
While condemning the Hamas attacks on innocent Israelis and the events of October 7, he said :"It was completely foreseeable that there would be action on the @ @ @ @ @ @ @ @ @ @ to the future, as things stand, it ' s only creating a recipe for future violence and hatred moving forward,"he said.
Reflecting on personal connections on all sides of the conflict, Hope expressed a desire to return to both Israeli and Palestinian friends with his family if the dust settles.
Comments : Our rules We want our comments to be a lively and valuable part of our community - a place where \textbf{readers} can debate and engage with the most important local issues.
The ability to comment on our stories is a privilege, not a right, however, and that privilege may be withdrawn if it is abused or misused. follow us This website and associated newspapers adhere to the Independent Press Standards Organisation ' s Editors ' Code of Practice.
If you have a complaint about the editorial content which relates to inaccuracy or intrusion, then \textbf{please} contact the editor here.
If you are dissatisfied with the response provided you can contact IPSO here

% matched lemmas: fundraising, please, reader, trend
\end{textsample}
